\chapter{Installation and configuration}
\label{ch:installation}

\section{Supported environment}

The package is developed and verified on Linux and Windows with CPython
3.11, on CPU.  The Tutorial (Chapter 2) gives the minimal recipe; this
chapter records the constraints behind it.

\begin{center}
\begin{tabular}{lll}
\toprule
Dependency & Constraint & Why it is pinned \\
\midrule
Python & $\geq 3.10$ & \texttt{int.bit\_count} in the determinant tools \\
numpy & $\geq 1.24$ & analysis layer, reference arithmetic \\
scipy & $\geq 1.10$ & L-BFGS-B and trust-region stages \\
pyscfad & $\geq 0.3.3$ & provides the traceable SCF and integral path \\
pyscf & $\geq 2.5$ & reference energies and comparisons \\
jax & $\geq 0.4$ & array layer; the exact pin is pyscfad's \\
matplotlib & any recent & figure sources only \\
\bottomrule
\end{tabular}
\end{center}

\texttt{pyscfad} pins a specific JAX version, which is why the package
should live in its own environment.  Mixing it into an environment with
a different JAX pin produces failures that look like numerical bugs and
are not.

\section{Runtime configuration}

Three settings are load-bearing.

\textbf{\texttt{jax\_enable\_x64 = True}.} The package is validated in
double precision.  With the default single precision the errors are not
loud: gradients lose several digits and some convergence gates become
unreachable, while energies look plausible.

\textbf{\texttt{JAX\_PLATFORMS = cpu}.} The CPU platform is the
verified configuration.  On a machine with a CUDA device, leaving the
platform unset selects the GPU, where the solver path used by the
orbital canonicalization raises an error unrelated to the chemistry.
Set the variable before importing JAX; setting it afterwards has no
effect.

\textbf{Thread environment variables.} The driver assumes it owns the
thread budget it was given.  Setting the five BLAS thread variables
keeps multiple concurrent calculations from oversubscribing a node;
their absence does not corrupt results, it makes timing meaningless and
slows everything.

\section{Verifying the installation}

\begin{lstlisting}[language=bash]
pytest tests/ -v
lint-imports        # or: pytest tests/test_layer_contracts.py
\end{lstlisting}

The suite is the primary artifact to look at: it contains the
first-order and second-order derivative checks against finite
differences, the cross-validation against PySCF, and the layer
contracts.  A passing suite on your machine is the intended proof that
the environment is sound.

\section{Optional: running the figures}

The figure sources under \texttt{figures/} import matplotlib and read
JSON results produced by the scripts; they require no additional
configuration.  They are not part of the package tests.

\section{Known environment friction}

\begin{itemize}
  \item The first evaluation of a freshly built node of the chain
    triggers XLA compilation that can take minutes.  This is expected;
    a stopwatch on a single evaluation will mislead.
  \item pyscfad's \texttt{Mole} objects cannot cross a
    \texttt{jax.jit} boundary.  The driver layer offers frozen-geometry
    factories for jitted contexts; the factory docstrings state which
    is which.
  \item On shared nodes, leave headroom for the owner of the other
    processes.  The memory profile of this package is dominated by
    compilation- and evaluation-time buffers of the differentiated
    chain, not by the integrals, and it grows steeply with the number
    of rotation parameters.
\end{itemize}
